\section{Related Work and Scope}
\label{sec:related}

\paragraph{Workspaces and agent-native file systems.}
OverlayFS provides layered namespaces with a private upper and shared
lower~\cite{overlayfs}. BranchFS applies copy-on-write branch contexts to
agentic exploration, with a commit/abort lifecycle described through a
workshop prototype and a proposed kernel abstraction~\cite{branchfs2026}.
YoloFS addresses visibility and control over agent effects through staging,
snapshots, and permissions~\cite{yolofs2026}; AgentFS documents a file-system
and tool-state abstraction for agents~\cite{agentfs2025}. These systems
establish workspace and agent-control mechanisms on which the present
problem builds. \sys focuses on serving private views through remote
metadata and object access. Process isolation, user approval, and the
commit semantics of other agent systems remain separate concerns.

Table~\ref{tab:comparison} identifies the comparison axes. Its last column
states the boundary to test against each baseline under matched persistence
and cache budgets; it reports no measured advantage.

\begin{table*}[t]
\centering
\scriptsize
\begin{tabularx}{\textwidth}{@{}p{1.60cm}p{1.95cm}p{2.15cm}p{2.0cm}Y@{}}
\toprule
System & Shared baseline & Private view & Remote metadata/data & Boundary tested in this work\\
\midrule
OverlayFS & local lower tree & kernel upper/copy-up & none by itself & remote immutable lower, explicit object reads, and durable upper\\
DeltaFS & immutable logs and manifests & per-job namespace & object store plus registry & current-value upper, directory-oriented pages, and online POSIX lookup\\
BranchFS & local COW branch context & branch commit/abort & prototype-local & remote persistence and multi-client lower sharing\\
CVMFS/Nydus & lazy immutable catalog/image & requires an external upper & cached objects/blobs & writable workspace semantics and metadata/data locator co-location\\
JuiceFS & chunk/slice/block volume & shared mutable POSIX view & object store plus KV & frozen metadata access, object formation at $g$, and persisted $(B,q)$ layout\\
\sys & native or packed baseline & optional private delta & objects plus metadata/control KV & separate and combined effects of view composition, packing, and object layout\\
\bottomrule
\end{tabularx}
\caption{Comparison axes used to choose baselines. Existing systems already
provide individual mechanisms; the opportunity examined here is their
independent use and optional composition for remote agent workloads.}
\label{tab:comparison}
\end{table*}

\paragraph{Namespace and metadata organization.}
TableFS, IndexFS, LocoFS, and InfiniFS organize, partition, and cache
metadata to change the work required by file-system
operations~\cite{tablefs2013,indexfs2014,locofs2017,infinifs2022}.
Recent systems also study hierarchical object metadata and distributed
update coordination~\cite{mantle2025,switchfs2026}. Their access paths make
record layout, batching, and cache reuse relevant comparison criteria:
a metadata service need not issue one remote request per file operation.

DeltaFS is a particularly close conceptual predecessor. It uses private
namespaces and immutable mutation logs over shared storage, with selective
composition of prior state~\cite{deltafs2021}. \sys builds on the separation
of shared state and private changes, and examines the metadata access units
used to serve a file view. Packed Metadata places common attributes and
data locators in indexed, bounded groups. It can serve a standalone
read-only view or the immutable lower of an online private workspace,
while the workspace contract also admits a native metadata lower. This
separation permits comparing the representation independently of view
composition, then measuring the cost of combining them.

\paragraph{Immutable distribution and storage backends.}
CernVM-FS provides demand-driven software distribution with indexed
catalogs, while Nydus describes metadata/data organization for lazy
image access~\cite{cvmfs2011,nydus_rafs}. Slacker and DADI also address
image distribution and startup costs~\cite{slacker2016,dadi2020}.
These systems supply close comparisons for on-demand access to an
immutable baseline. The distinguishing questions concern how records and
data locators are obtained, which fetched units can be reused, and how
construction and staging costs are amortized. The planned comparisons
cover both standalone read-only access and reuse beneath writable views.

Distributed systems such as Ceph and Tectonic separate aspects of metadata
and data services~\cite{ceph2006,tectonic2021}. Erasure-coded storage research considers
repair cost, changing redundancy over a file's lifetime, and the relation
between data placement and redundancy~\cite{azureec2012,morph2024,okapi2025}.
Okapi's separation of striping and redundancy grouping is different from
our $(B,q)$ distinction, although both motivate making coupled decisions
explicit. Client block granularity $g$ additionally determines how file
data forms objects; it is distinct from backend coding and physical
verification granularity. Our comparisons hold the code rate fixed and
separate object counts, checksums, tails, and repair work. Reducing a block
size can change these costs without changing the nominal redundancy ratio.
The same distinction applies to ordinary volumes, workspace data, and
packed snapshots where the selected object path supports layout selection.

\paragraph{Scope and unresolved evidence.}
The design targets three conditions: sufficient unchanged state to reuse
a baseline, sufficient reuse of immutable records to repay packing, and
an access pattern for which a different object or backend unit reduces
total storage work. These conditions need not occur together. Poor name
locality, high mutation, or rapid snapshot churn may favor a native
metadata path even when a different object layout is useful. The current
entry points and lifecycle support further bound the combinations that can
be tested (Section~\ref{sec:implementation}). Establishing each mechanism's
useful region and any additional value from composition remains the
purpose of the evaluation plan.
